This paper presented \repo{} as a software-and-methods workflow for nonlinear
finite-element solves and optimization with a maintained \jaxpetsc{} production
path. The central contribution is not a new constitutive theory or a universal
differentiable-PDE abstraction, but a maintained implementation pattern in
which local JAX differentiation, sparse PETSc nonlinear solves, and multiple
second-order routes are kept together across a heterogeneous benchmark suite.

The Plasticity3D evidence is intentionally narrower than a broad mechanics
claim. Same-case source-faithfulness to the source-style surrogate
implementation is strong, while the fixed-lambda source-operator comparison
does not support constitutive-equivalence language. The correct reading is
therefore that the repository provides a supported maintained surrogate
workflow for this benchmark family, with constitutive AD emerging as the
preferred cost-quality route on the locked flagship case. On Hyperelasticity,
the matched companion baseline against JAX-FEM shows that the maintained code
can agree very closely with a recent external implementation when the problem
contract is made genuinely symmetric. The promoted historical Plasticity3D
scaling campaign and the distributed topology workflow then show that the same
codebase remains operational beyond small matched comparisons.

Future work should follow the boundaries made visible by these results:
\begin{itemize}
  \item extend the validation surface for path-dependent mechanics beyond the
        present surrogate-versus-source-operator comparison, including a true
        incremental-history reference;
  \item add more fairness-first external baselines on tightly matched problem
        contracts;
  \item continue reducing linear-setup and multigrid costs in the flagship 3D
        mechanics workflows.
\end{itemize}

The next scientific step is not to overclaim this implementation, but
to turn this maintained workflow into a stronger reusable template for
additional nonlinear multiphysics and design problems with equally explicit
validation and fairness criteria.
